\documentclass[12pt,a4paper]{article}

\usepackage[T1]{fontenc}
\usepackage[utf8]{inputenc}
\usepackage[french]{babel}
\usepackage[margin=2.5cm]{geometry}
\usepackage{hyperref}
\usepackage{xcolor}
\usepackage{tcolorbox}
\usepackage{tikz}
\usepackage{tikz-qtree}
\usepackage{booktabs}
\usepackage{enumitem}
\usepackage{lmodern}
\usepackage{microtype}
\usepackage{titlesec}
\usepackage{fancyhdr}
\usepackage{graphicx}

\usetikzlibrary{shapes.geometric, arrows.meta, positioning, fit, backgrounds, calc}

% Colors
\definecolor{agentblue}{RGB}{30, 90, 160}
\definecolor{agentgreen}{RGB}{34, 139, 34}
\definecolor{agentred}{RGB}{180, 30, 30}
\definecolor{agentorange}{RGB}{210, 120, 0}
\definecolor{agentgray}{RGB}{240, 240, 245}
\definecolor{agentlightblue}{RGB}{220, 235, 255}
\definecolor{agentlightgreen}{RGB}{220, 245, 220}
\definecolor{agentlightred}{RGB}{255, 230, 230}

% tcolorbox styles
\tcbuselibrary{skins, breakable}

\newtcolorbox{capabilitybox}[1][]{
  colback=agentlightgreen,
  colframe=agentgreen,
  fonttitle=\bfseries,
  title=#1,
  breakable,
  arc=4pt
}

\newtcolorbox{limitbox}[1][]{
  colback=agentlightred,
  colframe=agentred,
  fonttitle=\bfseries,
  title=#1,
  breakable,
  arc=4pt
}

\newtcolorbox{improvebox}[1][]{
  colback=agentlightblue,
  colframe=agentblue,
  fonttitle=\bfseries,
  title=#1,
  breakable,
  arc=4pt
}

% Header/footer
\pagestyle{fancy}
\fancyhf{}
\fancyhead[L]{\textcolor{agentblue}{\textbf{AgentActuariat}}}
\fancyhead[R]{\textcolor{gray}{\small Analyse des capacités}}
\fancyfoot[C]{\thepage}
\renewcommand{\headrulewidth}{0.4pt}

% Section formatting
\titleformat{\section}{\large\bfseries\color{agentblue}}{}{0em}{}[\titlerule]
\titleformat{\subsection}{\normalsize\bfseries\color{agentblue!70!black}}{}{0em}{}

\hypersetup{
  colorlinks=true,
  linkcolor=agentblue,
  urlcolor=agentblue,
  pdftitle={AgentActuariat — Capacités et limites},
  pdfauthor={AgentActuariat}
}

% ─────────────────────────────────────────────
\begin{document}

\begin{titlepage}
  \centering
  \vspace*{2cm}
  {\Huge\bfseries\color{agentblue} AgentActuariat\par}
  \vspace{0.5cm}
  {\large\color{gray} Analyse des capacités et des limites du système\par}
  \vspace{2cm}
  \rule{0.6\textwidth}{0.5pt}
  \vspace{1cm}

  \begin{tcolorbox}[colback=agentgray, colframe=agentblue!30, width=0.75\textwidth, arc=6pt]
    \centering
    \textbf{Objet :} Ce document décrit précisément ce que fait l'agent,
    ce qu'il ne fait pas, et les priorités d'amélioration identifiées
    à partir de l'analyse de son architecture (avril 2026).
  \end{tcolorbox}

  \vfill
  {\small\color{gray} Document généré à partir de \texttt{docs/AGENT\_ARCHITECTURE.md}\par}
  {\small\color{gray} Avril 2026\par}
\end{titlepage}

\tableofcontents
\newpage

% ─────────────────────────────────────────────
\section{Vue d'ensemble — Architecture du système}

\texttt{AgentActuariat} est un agent multi-nœud LangGraph dont l'objectif est
de construire et certifier des tables de mortalité d'expérience à partir d'un
portefeuille assurantiel, puis de générer un rapport actuariel professionnel au
format PDF.

\bigskip

% ── TikZ diagram ──────────────────────────────
\subsection{Schéma de la pipeline}

\begin{figure}[h!]
\centering
\resizebox{\textwidth}{!}{%
\begin{tikzpicture}[
  node distance = 0.9cm and 1.4cm,
  box/.style = {rectangle, rounded corners=4pt, minimum width=2.6cm,
                minimum height=0.9cm, text centered, font=\small\bfseries,
                draw, line width=0.8pt},
  input/.style  = {box, fill=agentgray,    draw=gray},
  master/.style = {box, fill=agentlightblue, draw=agentblue},
  builder/.style= {box, fill=agentlightgreen,draw=agentgreen},
  writer/.style = {box, fill=agentlightblue, draw=agentblue},
  pipeline/.style={box, fill=agentlightblue!60, draw=agentblue!60,
                   minimum width=2.2cm, minimum height=0.75cm, font=\footnotesize},
  output/.style = {box, fill=agentlightred, draw=agentred},
  tools/.style  = {box, fill=agentlightgreen!60, draw=agentgreen!70,
                   minimum width=2.2cm, minimum height=0.75cm, font=\footnotesize},
  arr/.style    = {-{Stealth[length=5pt]}, line width=0.7pt},
  arrblue/.style= {arr, color=agentblue},
  arrgreen/.style={arr, color=agentgreen},
  arrred/.style = {arr, color=agentred},
  label/.style  = {font=\tiny\itshape, text=gray}
]

% ── Nodes ──────────────────────────────────────
\node[input]   (csv)     {CSV\\ Portefeuille};
\node[master,  right=of csv]  (master)  {MasterAgent\\ \footnotesize(intent + routing)};
\node[builder, right=of master](builder){BuilderAgent\\ \footnotesize(ReAct, max 5 iter.)};
\node[tools,   below=0.7cm of builder]  (tools)   {Tools actuariels\\ \footnotesize(exposure, lissage…)};
\node[writer,  right=of builder](writer){WriterNode\\ \footnotesize(no LLM)};

% pipeline steps to the right of writer, stacked vertically
\node[pipeline, right=1.6cm of writer, yshift=2.2cm] (p01) {01 load\_plan};
\node[pipeline, below=0.3cm of p01] (p02) {02 validation\_plan\\ \footnotesize(LLM)};
\node[pipeline, below=0.3cm of p02] (p03) {03 completion\_plan\\ \footnotesize(RAG ChromaDB)};
\node[pipeline, below=0.3cm of p03] (p04) {04 redaction\\ \footnotesize(LLM × section)};
\node[pipeline, below=0.3cm of p04] (p05) {05 assemble\\ \footnotesize(ReportLab)};
\node[pipeline, below=0.3cm of p05] (p06) {06 validate\_report\\ \footnotesize(LLM, verdict)};
\node[output,   below=0.5cm of p06] (pdf) {PDF\\ Rapport final};

% ── Arrows ─────────────────────────────────────
\draw[arrblue]  (csv)     -- node[above, label]{upload} (master);
\draw[arrblue]  (master)  -- node[above, label]{build} (builder);
\draw[arrgreen] (builder) -- node[right, label]{tool calls} (tools);
\draw[arrgreen] (tools)   -- ++(0,0.7) -| (builder);   % loop back
\draw[arrblue]  (builder) -- node[above, label]{BUILD\_DONE} (writer);

% master ↔ builder feedback
\draw[arrblue, dashed, bend right=30]
  (builder.south west) to node[below, label]{to\_master} (master.south east);

% writer → pipeline
\draw[arrblue]  (writer.east) -- ++(0.5,0) |- (p01.west);

% pipeline steps
\draw[arrblue] (p01) -- (p02);
\draw[arrblue] (p02) -- (p03);
\draw[arrblue] (p03) -- (p04);
\draw[arrblue] (p04) -- (p05);
\draw[arrblue] (p05) -- (p06);
\draw[arrred]  (p06) -- (pdf);

% need_data feedback
\draw[arrred, dashed, bend left=50]
  (p02.west) to node[left, label]{need\_data} (master.north east);

% retry step 04
\draw[arrblue, dashed, bend left=40]
  (p06.east) to[out=0, in=0] node[right, label]{retry minor} (p04.east);

% ── Legend box ─────────────────────────────────
\begin{scope}[shift={(-0.2,-3.8)}]
  \node[input,  minimum width=1.8cm, minimum height=0.55cm, font=\tiny] (l1) at (0,0) {Données};
  \node[master, minimum width=1.8cm, minimum height=0.55cm, font=\tiny] (l2) at (2.2,0) {Nœud LangGraph};
  \node[builder,minimum width=1.8cm, minimum height=0.55cm, font=\tiny] (l3) at (4.4,0) {Calcul / LLM actif};
  \node[output, minimum width=1.8cm, minimum height=0.55cm, font=\tiny] (l4) at (6.6,0) {Sortie};
  \node[font=\tiny\itshape, text=gray] at (3.3,-0.6) {Légende};
\end{scope}

\end{tikzpicture}
}
\caption{Architecture complète d'AgentActuariat — pipeline de bout en bout}
\label{fig:architecture}
\end{figure}

\subsection{Composants et modèle LLM}

\begin{center}
\begin{tabular}{llll}
\toprule
\textbf{Composant} & \textbf{Rôle} & \textbf{LLM} & \textbf{Type} \\
\midrule
MasterAgent    & Routing, intent, désambiguïsation & GPT-4o & 1–3 appels/tour \\
BuilderAgent   & Calculs actuariels (ReAct)        & GPT-4o & Borné à 5 iter. \\
WriterNode     & Déclencheur pipeline              & Aucun  & Déterministe \\
Étape 02       & Validation données suffisantes    & GPT-4o & JSON mode \\
Étape 04       & Rédaction narrative               & GPT-4o & temp=0.4 \\
Étape 06       & Contrôle qualité final            & GPT-4o & JSON mode \\
\bottomrule
\end{tabular}
\end{center}

\newpage

% ─────────────────────────────────────────────
\section{Ce que l'agent fait}

\begin{capabilitybox}[Capacités validées par l'analyse de l'architecture]

\subsection*{1. Ingestion et préparation des données}
\begin{itemize}[leftmargin=1.5em]
  \item Charge un fichier CSV de portefeuille assurantiel via l'interface Dash.
  \item Mappe automatiquement les colonnes sur les champs attendus
        (\texttt{date\_naissance}, \texttt{date\_entree}, \texttt{date\_sortie},
        \texttt{cause\_sortie}) par règles d'abord, puis par LLM si ambigu.
  \item Désambiguïse le plan d'étude au premier tour (objectif, filtres, tranches d'âge).
  \item Persiste le DataFrame en Parquet pour éviter de le retransmettre à chaque tour.
\end{itemize}

\subsection*{2. Calcul actuariel complet}
Le BuilderAgent orchestre via function-calling l'enchaînement des outils suivants :
\begin{itemize}[leftmargin=1.5em]
  \item \textbf{Exposition} : calcul en années-personnes par âge.
  \item \textbf{Taux bruts} ($\hat{q}_x$) : méthode Kaplan-Meier ou actuarielle classique.
  \item \textbf{Lissage} ($q_x$ lissés) : ajustement paramétrique avec contrôle de monotonie.
  \item \textbf{Diagnostics} : détection des âges à faible crédibilité, recommandation de régime.
  \item \textbf{Validation statistique} : intervalles de confiance, test du $\chi^2$.
  \item \textbf{Benchmarking} : abattement vs table réglementaire, SMR global.
  \item \textbf{Régression de Cox} : rapport de risque par sexe avec IC à 95\%.
  \item \textbf{Régression logit} : ajustement $q_x$ lissés vs table de référence ($R^2$, pente).
  \item \textbf{Comparaison antérieure} : dérive globale et âges à forte dérive vs étude précédente.
\end{itemize}

\subsection*{3. Génération du rapport PDF}
Pipeline de 6 étapes séquentielles :
\begin{enumerate}[leftmargin=1.5em]
  \item \textbf{Chargement du plan} : résolution des \texttt{\{\{\ placeholders\ \}\}} depuis le
        \texttt{data\_store}, construction d'un prompt autonome par section.
  \item \textbf{Validation des données} : vérification LLM que chaque section a ses inputs
        — retour au builder si manquant.
  \item \textbf{Complétion RAG} : injection d'exemples de rédaction depuis ChromaDB
        (rapport de référence indexé).
  \item \textbf{Rédaction} : LLM GPT-4o rédige chaque section en parallèle
        (\texttt{temperature=0.4}, \texttt{max\_tokens=1200}).
        Tables et graphiques sont rendus de façon déterministe (ReportLab / matplotlib)
        \emph{avant} l'appel LLM.
  \item \textbf{Assemblage} : ReportLab assemble les sections dans l'ordre fixe.
  \item \textbf{Validation finale} : LLM audite le rapport (verdict ok / minor / major)
        avec retry ciblé sur les sections KO.
\end{enumerate}

\subsection*{4. Validator de traçabilité des chiffres}
\begin{itemize}[leftmargin=1.5em]
  \item Extrait tous les nombres cités dans le texte narratif par regex.
  \item Vérifie que chaque nombre est traçable dans le \texttt{data\_store}
        (tolérance $\pm 2\,\%$, conversions fraction/pourcent/millième).
  \item En cas d'échec : retry ciblé avec feedback chirurgical au LLM (1 fois max).
  \item Whitelist pour constantes statistiques universelles (1,96 ; 0,05 ; etc.).
\end{itemize}

\subsection*{5. Mémoire persistante multi-couches}
\begin{itemize}[leftmargin=1.5em]
  \item \textbf{Working memory} (RAM) : état LangGraph complet, compacté au-delà de 15 messages.
  \item \textbf{Business memory} (JSON sur disque) : plan d'étude, mapping colonnes, résultats
        des outils actuariels — survit au redémarrage.
  \item \textbf{Dataset memory} (Parquet sur disque) : DataFrame brut — écrit une seule fois.
  \item \textbf{Conversation memory} : historique LangGraph avec résumé structuré automatique.
\end{itemize}

\end{capabilitybox}

\newpage

% ─────────────────────────────────────────────
\section{Ce que l'agent ne fait pas — Limites identifiées}

\begin{limitbox}[Limites confirmées par l'analyse du code]

\subsection*{1. Sections narratives potentiellement tronquées}
\texttt{max\_tokens=1200} par section à l'étape 04. Une section actuarielle professionnelle
(construction de table, analyse statistique) peut nécessiter 600–900 mots, soit
environ 1\,200–1\,500 tokens. Ce plafond est donc \emph{au bord de l'insuffisance}
pour les sections les plus denses.

\subsection*{2. Température trop élevée sur les étapes critiques}
Les étapes 02 (validation plan) et 06 (contrôle qualité) n'ont pas de
\texttt{temperature} explicite dans le code $\Rightarrow$ défaut OpenAI = \textbf{1,0}.
Pour un contenu de certification réglementaire, une température de 1,0 introduit
une variabilité inacceptable dans les verdicts JSON.

\subsection*{3. RAG insuffisant pour capturer le style de référence}
L'étape 03 injecte au maximum \textbf{600 caractères} par section depuis le rapport
d'exemple (ChromaDB). C'est trop peu pour transmettre le registre, la structure
et le niveau de détail d'un rapport actuariel professionnel.

\subsection*{4. WriterNode sans intelligence contextuelle}
Le WriterNode appelle \texttt{run\_pipeline()} directement \emph{sans aucun appel LLM
au niveau du nœud}. Il n'y a donc aucune adaptation intelligente à la spécificité
des données (portefeuille atypique, âges extrêmes, crédibilité faible) avant
de lancer la rédaction.

\subsection*{5. Ordre d'exécution des outils non garanti}
Le BuilderAgent choisit librement l'ordre des tools via function-calling.
Les garde-fous sont \emph{passifs} uniquement : un tool retourne une erreur si son
prérequis manque, mais l'agent peut théoriquement appeler \texttt{benchmarking}
avant \texttt{smoothing} sans blocage proactif.

\subsection*{6. Working memory perdue au redémarrage}
Le \texttt{MemorySaver} de LangGraph est en RAM. Un redémarrage du process efface
l'état LangGraph de toutes les sessions actives. Seule la \texttt{SessionState}
(JSON/Parquet) survit.

\subsection*{7. Contenu des system prompts non auditable en l'état}
Les prompts du MasterAgent et du BuilderAgent sont chargés dynamiquement depuis
des fichiers \texttt{.md} dans \texttt{agents/*/agent\_instructions/}. Leur contenu
exact n'a pas été vérifié dans cette analyse — ils constituent un angle mort
de la garantie qualité.

\subsection*{8. Pas de gestion robuste des CSV mal formés}
La désambiguïsation peut boucler ou échouer silencieusement sur un CSV mal structuré
(encodage, colonnes manquantes, formats de date hétérogènes). Aucun mécanisme
d'abandon propre n'a été identifié dans le code.

\end{limitbox}

\newpage

% ─────────────────────────────────────────────
\section{Points d'amélioration prioritaires}

\begin{improvebox}[Recommandations classées par impact / effort]

\subsection*{Priorité 1 — Corriger la température des étapes 02 et 06}
\textbf{Impact : élevé — Effort : minimal (2 lignes de code)}\\
Fixer \texttt{temperature=0.0} ou \texttt{0.1} sur les appels LLM des étapes de validation
(02 et 06). Un verdict de certification doit être déterministe.

\begin{center}
\begin{tcolorbox}[colback=black!5, colframe=black!30, width=0.8\textwidth,
                  fontupper=\ttfamily\small, arc=3pt]
\# Étape 02 et 06 — ajouter :\\
temperature=0.1,~~\# déterminisme pour contenu réglementaire
\end{tcolorbox}
\end{center}

\subsection*{Priorité 2 — Augmenter max\_tokens à l'étape 04}
\textbf{Impact : élevé — Effort : faible}\\
Passer \texttt{max\_tokens} de 1\,200 à \textbf{2\,000} pour les sections denses
(\texttt{construction}, \texttt{analysis}), en conservant 1\,200 pour les sections
courtes (\texttt{preamble}, \texttt{conclusion}).

\subsection*{Priorité 3 — Enrichir le RAG de l'étape 03}
\textbf{Impact : élevé — Effort : moyen}\\
Porter la limite de troncature de 600 à \textbf{1\,500 caractères} par section,
et augmenter le nombre de chunks récupérés de 1 à 3. L'exemple de référence
doit couvrir au moins un paragraphe complet pour transmettre le style.

\subsection*{Priorité 4 — Introduire un LLM de contexte dans WriterNode}
\textbf{Impact : moyen — Effort : moyen}\\
Avant de déclencher \texttt{run\_pipeline()}, ajouter un appel LLM léger
(\texttt{gpt-4o-mini}, \texttt{max\_tokens=300}) qui analyse le \texttt{data\_store}
et produit un \emph{briefing contextuel} injecté dans les prompts de section :
particularités du portefeuille, âges à crédibilité faible, anomalies détectées.

\subsection*{Priorité 5 — Garde-fous proactifs sur l'ordre des tools}
\textbf{Impact : moyen — Effort : moyen}\\
Ajouter une vérification explicite dans \texttt{tools\_node.py} :
refuser l'exécution de \texttt{benchmarking} ou \texttt{logit\_regression}
si les prérequis (\texttt{smoothed\_table}) sont absents du \texttt{data\_store},
avec un message d'erreur clair renvoyé au BuilderAgent.

\end{improvebox}

\bigskip

% ─────────────────────────────────────────────
\section*{Synthèse}

\begin{center}
\begin{tabular}{p{6cm} p{6cm}}
\toprule
\textbf{\color{agentgreen}Points forts} & \textbf{\color{agentred}Points faibles} \\
\midrule
Pipeline déterministe (tables, graphes) & max\_tokens trop bas pour sections denses \\
Validator de traçabilité des chiffres   & Temperature par défaut = 1.0 (étapes 02/06) \\
Mémoire persistante multi-couches       & RAG insuffisant (600 chars) \\
Retry ciblé sur sections KO             & WriterNode sans intelligence contextuelle \\
Bornage du ReAct à 5 itérations         & Ordre des tools non garanti (passif) \\
\bottomrule
\end{tabular}
\end{center}

\vfill
\begin{center}
  \textcolor{gray}{\small Document produit à partir de \texttt{docs/AGENT\_ARCHITECTURE.md} — Avril 2026}
\end{center}

\end{document}
